% from ega4-16-fr.tex
\oldpage[IV-4]{48}
  prenant $X=\operatorname{Spec}(A)$, où $A$ est un anneau de Dedekind, et
  $\mathscr{I}=\widetilde{\mathfrak{I}}$, où $\mathfrak{I}$ est un idéal
  premier non principal de $A$; on a alors en effet
  $\mathscr{I}_y/\mathscr{I}_y^2=0$ en tout point $y$ distinct du point
  $x\in X$ correspondant à $\mathfrak{I}$, et
  $\mathscr{I}_x/\mathscr{I}_x^2$ est de rang $1$ sur le corps
  $\mathcal{O}_x/\mathscr{I}_x$; en outre $\mathscr{I}$ est évidemment un
  Idéal régulier.
  \item[(ii)] Dans (16.9.5), on ne peut remplacer « quasi-régulière » par
  « régulière », même lorsque $X$ est un préschéma (cf. (16.9.12)).
  Désignons en effet par $B$ l'anneau des germes de fonctions indéfiniment
  différentiables au point $0$ dans $\mathbb{R}$; il a un idéal maximal
  $\mathfrak{m}$ engendré par le germe $t$ de l'application identique de
  $\mathbb{R}$ au point $0$, et l'intersection $\mathfrak{n}$ des
  $\mathfrak{m}^k$ pour $k>0$ n'est pas réduite à $0$. Soit maintenant $A$
  l'anneau quotient $B[T]/\mathfrak{n}TB[T]$, et soient $f_1,f_2$ les images
  canoniques dans $A$ des éléments $t$ et $T$ de $B[T]$. La suite
  $(f_1,f_2)$ est régulière dans $A$: en effet, $f_1$ n'est pas diviseur de
  $0$ dans $A$, car la relation $tP[T]\in\mathfrak{n}TB[T]$, pour un
  polynôme $P\in B[T]$, entraîne que les produits de $t$ par les coefficients
  de $P$ appartiennent à l'idéal $\mathfrak{n}$, et il en résulte aussitôt que
  ces coefficients sont eux-mêmes dans $\mathfrak{n}$, donc
  $P[T]\in\mathfrak{n}TB[T]$. Comme $B/tB$ est isomorphe à $\mathbb{R}$,
  $A/f_1A$ est isomorphe à l'anneau de polynômes $\mathbb{R}[T]$, donc
  intègre, et l'image de $f_2$ dans $A/f_1A$, étant égale à $T$, n'est donc
  pas diviseur de $0$, ce qui prouve notre assertion. Cependant $f_2$ est
  diviseur de $0$ dans $A$, car pour tout élément $x\in\mathfrak{n}$ non nul,
  l'image de $x$ dans $A$ est $\ne0$, mais l'image de $xT$ est nulle. On en
  conclut que la suite $(f_2,f_1)$ n'est pas régulière dans $A$; d'autre part
  l'idéal $\mathfrak{I}=f_1A+f_2A$ est distinct de $A$, donc les conditions
  b), b') et c) de (16.9.5) n'entraînent pas la condition a) lorsqu'on y
  remplace « quasi-régulière » par « régulière ».
\end{enumerate}
\end{remarks}

\begin{env}[16.9.7]
\EGAReaderTarget{ega4.c04.statement.16-9-7}{IV.16.9.7}
Si $X=\operatorname{Spec}(A)$ est un schéma affine, nous dirons qu'un idéal
$\mathfrak{I}$ de $A$ est régulier (resp. quasi-régulier) si l'Idéal
$\mathscr{I}=\widetilde{\mathfrak{I}}$ de $\mathcal{O}_X$ est régulier
(resp. quasi-régulier); on notera que cette notion est locale et n'implique
nullement l'existence d'un système de générateurs de $\mathfrak{I}$ formant
dans $A$ une suite régulière (resp. quasi-régulière) comme le montre l'exemple
(16.9.6, (i)); il en est toutefois ainsi lorsque $A$ est local (16.9.5).

La proposition (16.9.4) se traduit en termes d'immersions quasi-régulières de
la façon suivante:
\end{env}

\begin{proposition}[16.9.8]
\label{IV.16.9.8}
Soit $j:Y\to X$ un morphisme de préschémas; pour que $j$ soit une immersion
quasi-régulière, il faut et il suffit que $j$ satisfasse aux conditions
suivantes:
\begin{enumerate}
  \item[(i)] $j$ est une immersion localement de présentation finie.
  \item[(ii)] Le faisceau conormal
  $\operatorname{Gr}^1(j)=\mathcal{N}_{Y/X}$ (16.1.2) est un
  $\mathcal{O}_Y$-Module localement libre.
  \item[(iii)] L'homomorphisme canonique
  \[
    \operatorname{S}_{\mathcal{O}_Y}^{\bullet}
      (\operatorname{Gr}^1(j))\longrightarrow\operatorname{Gr}^{\bullet}(j)
  \]
  (16.1.2.2) est bijectif.
\end{enumerate}
\end{proposition}

La question étant locale sur $Y$, on peut se borner au cas où $j$ est
l'injection canonique d'un sous-préschéma fermé $Y$ de $X$, et alors la
traduction de (16.9.4) en (16.9.8) résulte de l'explicitation de
$\operatorname{Gr}^1(j)$ et $\operatorname{Gr}^{\bullet}(j)$ en termes de
l'Idéal $\mathscr{I}$ de $\mathcal{O}_X$ définissant le sous-préschéma $Y$
(16.1.3, (ii)).

\begin{corollary}[16.9.9]
\label{IV.16.9.9}
Soient $Y$ un préschéma, $X$ un $Y$-préschéma, $j:Y\to X$ une $Y$-section
de $X$, de sorte que le $n$-ème invariant normal $\mathscr{A}^{(n)}$ de $j$
(16.1.2) est une $\mathcal{O}_Y$-Algèbre augmentée (16.1.7); posons
$\mathscr{A}^{(\infty)}=\varprojlim\mathscr{A}^{(n)}$. Pour que $j$ soit une
immersion quasi-régulière, il faut et il suffit que $j$ soit localement de
présentation finie, et que tout $y\in Y$ admette un voisinage ouvert affine
$U$ d'anneau $C$ tel que $\mathscr{A}^{(\infty)}|U$ soit isomorphe comme
$\mathcal{O}_U$-Algèbre topologique augmentée à
$\mathcal{O}_U[[T_1,\ldots,T_n]]$.
\end{corollary}

On peut se borner au cas où $j$ est une immersion fermée en se restreignant
à un voisinage assez petit de $y$ (voir le raisonnement de (16.4.11)), et
alors $\mathcal{O}_Y$ s'identifie à une Algèbre quotient
$\mathcal{O}_X/\mathscr{I}$ et l'homomorphisme surjectif canonique
$\mathcal{O}_X\to\mathcal{O}_Y$ admet un
